\section{总结与延伸}

\subsection{本讲核心要点}

\begin{enumerate}
    \item \textbf{图像是数字}：计算机通过 $H \times W \times C$ 的数字矩阵来表示图像，其中 $C$ 是通道数（灰度图 $C=1$，彩色图 $C=3$）。

    \item \textbf{特征提取是视觉识别的核心}：要识别物体，需要提取能区分不同类别的特征。手工特征设计面临递归性和变异性的困境，深度学习通过从数据中自动学习特征层次结构来解决这一问题。

    \item \textbf{全连接网络不适合图像}：展平操作丧失空间信息，且参数量过大。

    \item \textbf{卷积操作的三大优势}：
    \begin{itemize}
        \item 局部连接——保留空间结构
        \item 参数共享——大幅减少参数量
        \item 平移等变——同一模式无论出现在图像哪个位置都能被检测到
    \end{itemize}

    \item \textbf{CNN 架构}：卷积层 + ReLU + 池化层的堆叠构成特征学习模块，全连接层 + Softmax 构成分类模块。

    \item \textbf{CNN 是通用视觉架构}：同样的特征学习框架可以用于分类、目标检测、语义分割等多种任务。
\end{enumerate}

\subsection{与前两讲的联系}

\begin{table}[H]
\centering
\begin{tabular}{p{3cm}p{5cm}p{5cm}}
\toprule
\textbf{方面} & \textbf{Lecture 1 (全连接网络)} & \textbf{Lecture 3 (CNN)} \\
\midrule
输入形式 & 一维特征向量 & 二维/三维图像张量 \\
连接方式 & 全连接（所有输入到所有神经元） & 局部连接（每个神经元只看局部区域） \\
参数共享 & 无 & 滤波器在空间上共享 \\
空间信息 & 不保留 & 保留并利用 \\
核心操作 & 矩阵乘法 & 卷积操作 \\
\bottomrule
\end{tabular}
\caption{全连接网络与 CNN 的对比}
\end{table}

\begin{importantbox}{核心哲学：让数据结构指导架构设计}
CNN 的设计哲学是\textbf{利用数据的内在结构来设计网络架构}。图像具有空间局部性——相邻像素之间的关系比远距像素更密切；图像具有平移不变性——同一模式可以出现在任何位置。卷积操作完美地利用了这两个先验知识。这一哲学在深度学习中反复出现：RNN 利用了序列的时间结构，Transformer 中的注意力机制利用了元素间的关系结构。
\end{importantbox}

\subsection{拓展阅读}

\begin{itemize}
    \item MIT 6.S191 课程官网：\url{http://introtodeeplearning.com}
    \item LeCun et al. (1998), "Gradient-Based Learning Applied to Document Recognition" — CNN 的奠基之作
    \item Krizhevsky et al. (2012), "ImageNet Classification with Deep Convolutional Neural Networks" (AlexNet) — 深度学习在计算机视觉领域的里程碑
    \item Lee et al. (2009), ICML — 本讲中特征层次结构可视化的原始论文
    \item McKinney et al. (2020), Nature — CNN 在乳腺癌筛查中超越人类医生的研究
    \item Esteva et al. (2017), Nature — CNN 用于皮肤癌诊断的研究
    \item Rohrer, "How do CNNs work?" — 本讲中 X 检测案例的参考
    \item Karn, "An Intuitive Explanation of CNNs" — 本讲中卷积操作动画的参考
\end{itemize}

%% --- Fin del contenido --- %%

\end{document}
